\chapter[L'usage au centre du changement]{L'usager au centre de ce changement : la traduction des besoins en services fonctionnels et ergonomiques}

\lettrine{L}a valeur scientifique d'une bibliothèque numérique se mesure aussi à la conception des services mis en place par la bibliothèque pour répondre aux besoins identifiés de ses usagers. Le passage au prototype Codex représente un changement majeur pour la bibliothèque Cujas, permettant de mettre ne place une éditorialisation, adaptée aux différents types de public du lecteur éclairé au chercheur chevronné.
Dès lors, les retours d'expérience des chercheurs interrogés révèlent l'attente fondamentale: celle de retrouver un confort d'étude et la maniabilité aussi proches que celle d'un livre physique. Le défi de la refonte de la bibliothèque numérique consiste à concilier deux attentes : la \enquote{trouvabilité} et la \enquote{découvrabilité}, s'adaptant ainsi aux pratiques contemporaines de la recherche.

\section{L'interface de recherche et d'affichage}
Historiquement, CujasNum révèle que si le moteur de recherche simple s'avérait robuste et tolérant aux troncatures, l'interface souffrait quant à elle de visibilité et d'une confusion entre les différents modes de recherche. Le prototype Codex a jeté les bases de la modernisation, mais révèle des limites ergonomiques en multipliant et en répétant les onglets de recherche à différents endroits, créant une surchage cognitive. Pour surmonter ces inconvénients et répondre aux attentes formulées par des chercheurs articulent les spécificités autour de trois modalités : la recherche épurée, des filtres restructurées et des options de tri.

\subsection{La recherche simple par barre unique}
Sur CujasNum, les usagers étaient régulièrement confrontés à un bruit documentaire, ou à l'inverse à des silences inexplicables causé par la confusion entre la recherche dans les notices et la recherche plein-texte. Pour y remédier, la future page d'accueil de la bibliothèque numérique proposera une barre de recherche simple unique, positionnée au centre de l'interface. L'usager doit pouvoir sélectionner explicitement son périmètre via trois options de recherche simple : \enquote{à travers les métadonnées, la recherche s'effectue dans les notices descriptives ; à travers le plein-texte, la recherche interroge les textes océrisés et la dernière option est la combinaison des métadonnées et du plein-texte, la recherche s'effectue alors en croisant les données textuelles et descriptives de la plateforme}.

Ce modèle de recherche s'inspire des pratiques de recherche et d'autres institutions patrimoniales, comme la BIS ou l'INHA, ont développé au sein de leur propre bibliothèque numérique. Cette répartition permet d'accompagner le lecteur éliminant ainsi toute ambiguïté. Cette barre simple disposera d'un système d'autocomplétion ou de saisie intuitive, guidant l'usager dans sa recherche.

\subsection{La recherche avancée multicritères}
Dans CujasNum, le formulaire de recherche avancée était synonyme de confusion pour les agents et pour les usagers, car elle se situait \enquote{sur la même page que la recherche en texte intégrale} et utilisait un vocabulaire très professionnel, comme l'intitulé \enquote{notice bibliographique} peu compréhensible pour des étudiants ou des chercheurs non-spécialistes du livre.

La prototype Codex proposera un bouton de recherche avancée unique, clairement visibilité et situé en-dessous de la barre de recherche simple. C'est une pratique de présentation devenue courante que l'on retrouve dans différentes bibliothèques numériques, comme chez Lucienne de l'ENS ou encore Cariatide de l'INHA. Ce mode de présentation épurée mais visible réduit la charge cognitive des usagers, reproduisant des habitudes de recehrche des grands moteurs de recherche, comme Google ou Firefox. En cliquant sur le bouton de recherche avancée, l'utilisateur se trouve rediriger vers le formulaire qu'il n'a plus qu'à remplir. Ce formulaire permet ainsi de préciser des zones de recherche acceptant différents champs de saisie combiné à des opérateurs booléens.

Cette structuration permet aux chercheurs de réaliser des requêtes complexes en croisant les sources disponibles dans la bibliothèque numérique. Ce formulaire permet de répondre à l'une des attentes des chercheurs qui apprécient \enquote{de mener leur recherche à partir d'un point précis et élargissant progressivement leur périmètre de recherche au fur et à mesure}.

\subsection{Le filtrage par facettes}
L'un des chantiers identifiés au cours de l'audit de CujasNum se révèle être la gestion des filtres et des facettes. L'indexation historique reposait sur l'indexation Hauville. Cette indexatioin a généré des facettes ultra-spécifiques, souvent reliées à un seul document, comme la facette \enquote{Droit romain - Loi des Douze Tables}. Ce niveau de granularité entraînait une importante frustration pour les usagers signalant \enquote{qu'à cause de ça il n'utilisait ces facettes} ne souhaitant pas passer à côté d'un document. De surcroît, des variations de casse ou de typologies créaient davantage de facettes.

Une restructuration intégrale de ce parcours est envisagée par la bibliothèque Cujas. Le filtrage sera regroupé dans une nouvelle rubrique intitulée \enquote{Explorer}. La bibliothèque numérique proposera un ensemble de facettes générales et normalisées applicables à plusieurs types de documents. Cette rerstructuration permettra ainsi aux usagers de mieux comprendre comment sont présentés les collections. En effet, comme le démontre Leblanc \enquote{les facettes reflètent la manière dont la bibliothèque numérique a pensé ses collections et aident ainsi l'utilisateur à apprendre empiriquement l'organisation des collections}. Ce système de facettes permet ainsi la bibliothèque Cujas de guider intellectuellement les usagers. Un bon système de guidage bien formé et pensé devient un service indispensable \enquote{et permettent d'éviter le phénomène de surabondance des choix}.

Le système de filtrage ou de tri est tout aussi important que les autres modes de recherche, et nécessite d'être bien pensé. La simplification de l'indexation et des différents modes de recherche évite la confusion de l'usager, complétant ainsi les autres services proposés par la bibliothèque numérique. 
\section{L'Espace personnel et les services des chercheurs}
La refonte de la bibliothèque numérique de Cujas doit s'affranchir d'une consultation passive de l'interface en un véritable espace de travail personnalisé tout en y adossant des outils d'analyse scientifique avancées.

\subsection{L'intégratiion du standard IIIF et de la visionneuse Mirador}
Sur CujasNum, à l'origine, la consultation des ouvrages reposait sur une visionneuse JPEG, de bonnes qualités. De plus, elle offrait également la possibilité de télécharger le document consulté sans format PDF sans offrir différentes options de téléchargements ou d'autres formats. Un des chercheurs expliquait qu'il était peu courant qu'une bibliothèque offre différents options de téléchargement, qu'il trouvait cela dommage car par moment il aurait bien aimé pouvoir choisir uniquement de télécharger les pages intéressantes pour son étude. Une doctorante mentionne le fait que ce serait bien d'ajouter en plus des options de téléchargement \enquote{la possibilité de conserver la page de titre de l'ouvrage}. De cette façon, l'usager se souvient d'où ces pages ont été extraites.

Pour répondre à ces attentes, le prototype Codex, développé par la bibliothèque Cujas, sépare la couche de diffusion de celle de conservation, grâce à l'implémentation d'un serveur IIIF. Ce standard permet la diffusion dynamique d'image accompagnée d'outils scientifiques, proposant la possibilité de réaliser des captures d'écran, des annotations.

L'implémentation de la visionneuse Mirador offre l'opportunité aux chercheurs de la comparaison des documents au sein d'un même espace de travail. Les chercheurs peuvent désormais consulter plusieurs ouvrages afin de mener des analyses comparatives. Un autre des avantages est la possibilité d'utiliser un zoom profond sur une image exposée en haute définition. Elle offre aussi des réglages de rotations d'images, de contrastes et de colorimétrie.

\subsection{L'espace personnel du chercheur}
L'enquête de terrain et les entretiens menés auprès des usagers révèlent leur souhait d'avoir leurs propres espaces de  travail pour mener à bien leurs recherches. Sous CujasNum, il avait pris l'habitude de télécharger chaque document ne pouvant les conserver une fois le navigateur fermé.

Pour ancrer durablement l'usager dans la bibliothèque numérique de Cujas, l'implémentation d'un espace personnel est une priorité fonctionnelle majeure. Cet espace proposera des services simples et intuitifs comme la gestion des favoris, l'archivage des requêtes, et la recherche ciblée dans les sélections. L'espace personnel offre aussi la possibilité de rédiger des notes de travail ainsi que des annotations personnelles. Comme l'a expliqué Elina Leblanc, l'espace personnel répond à un besoin d'archivage d'usage. Il s'agit comme \enquote{une barrière à l'oubli} et génère \enquote{une économie de temps} cruciale pour les chercheurs.

\subsection{L'exposition des données}
Afin d'accompagner les usagers dans leurs rédactions, le prototype Codex doit offrir la possibilité des services d'exportation des données et d'analyse linguistiques aux standards du web sémantique.

Les chercheurs interrogés sont pour l'ajout de l'export de références bibliographiques normalisées, comme Zotero, leur permettant ainsi de gagner du temps. comme le souligne Antoine Courtin, \enquote{que les pratiques des chercheurs en termes de données, se limitent bien souvent à copier des données dans un document de traitemenet}. De plus, grâce à l'implémentation du module open-source \textit{Bibliography} d'Omeka S, la plateforme rendra possible l'exportation instantané de la notice bibliographique vers des gestionnaires de références, comme Zotero ou EndNote, sous des formats interopérables standardisés tels que BibTex, RIS ou JSON-LD.

L'intégration du format XML-TEI s'impose comme le standard pour encader la structure logique des textes juridiques, facilitant ainsi la reconnaissance des entités nommées (personnalités citées, lieux, date) et permet la fouille de texte. Comme l'expliquent Emmanuelle Bermès et Gautier Poupeau, la fouille de texte est rendue possible grâce à la disponibilité des corpus. \enquote{La \enquote{fouille de textes et de données}(en anglais \textit{text and data mining} ou TDM) embrasse un nouveau paradigme dans l'accès aux documents ou aux sources, qu'on pourrait désigner sous le nom de lecture distante, dans la lignée de Franco Moretti}.L'intérêt fondamental du TDM permet d'exploiter la masse de textes désormais numérique s'affranchissant des limites physiques de la lecture traditionnelle.

En dotant l'usager d'outil de pointe comme une visionneuse Mirador, un compte personnel et des formats structurés, la bibliothèque numérique devient un véritable outil pour mener à bien les projets en humanités numériques. Afin de mettre en avant cette collaboration entre les chercheurs et la bibliothèque, l'infrastructure doit alors s'accompagner d'une médiation intellectuelle. 
\section{Valorisation scientifique des collections}
L'éditorialisation constitue désormais un maillon indispensable pour donner de la cohérence intellectuelle à une bibliothèque numérique. L'éditorialisation contextualise les fonds, guide l'usager et valorise l'exxpertise scientifique de l'établissement.

Lors des enquêtes de terrains, les chercheurs ont fortement insisté sur la mise en place de cette médiation. L'un des chercheurs souligne qu'actuellement dans CujasNum, cela manque parfois de textes explicatifs, notamment sur la constitution des collections ou explicitant les choix de classifications. De même, un autre professeur rappelle l'importance d'attirer l'attention des usagers sur le contenu et la valeur singulière des ouvrages anciens.

Pour mettre en place cette médiation culturelle sans alourdir l'infrastructure technique, la bibliothèque Cujas peut s'inspirer du fonctionnement de la BIS. Lors de la conception d'expositions virtuelles ou de billets d'actualité, la BIS renonce à l'utilisation d'exposition virtuelles intégrées dans des visionneuses lourdes et peu portables. Elle a donc décidé d'opter pour le format Markdown. Les billets de blog et les pages d'expositions deviennent de simples fichiers textes structurés par des balises Markdown très légères, et exportables d'une solution technique à une autre. Ce choix technique permet de s'affranchir des instabilités des extensions de CMS, assurant une pérennité des données.

Pour rationaliser et homogénéiser la production éditoriale scientifique écrite par des tiers, la plateforme doit appliquer une démarche de qualité rigoureuse inspirée du guide de NuBIS. ce guide impose des règles strictes autour de trois points : la gestion des images, les contraintes rédactionnelles et le respect des droits d'auteur. Le respect de ces trois règles imposent un cadre stricte que doivent suivre les chercheurs souhaitant participer à l'éditorialisation de leur bibliothèque numérique. Comme l'a expliqué Elina Leblanc qui a analysé les \enquote{expositions nativement numériques} comme des espaces collaboratifs de recherche. Elle parle des avantages de permettre à des chercheurs, des enseignants ou des doctorants de rédiger des focus ou des notices scientifiques \enquote{participe au processus d'éditorialisation et conduit ainsi à créer de nouveaux savoirs, pouvant à leur tout être partagés}.

\subsection{Les formats de valorisation}
En s'appuyant sur ces règles de rédaction, l'éditorialisation scientifique de la BIU Cujas se déploiera sous plusieurs formes complémentaires, penser pour différents niveaux de lecture, comme la Galerie des grands juristes et l'éditorialisation par des juristes pour des juristes.

La \enquote{Galerie des grands juristes} est proposé par le directeur de la bibliothèque Cujas. Cette entrée thématique prendra la forme d'une galerie de portraits intéractifs. Inspirée de la même idéée que la bibliothèque de l'ENS qui a mis en avant ces collections, sous le nom des \textit{Parcours normalien.ne.s}. Chaque juriste fera l'objet d'une notice biographique, compris entre 5 000 et 10 000 signes, rédigée par un chercheur. Depuis la fiche bibliographique de l'auteur, l'usager pourra rebondir instantanément sur l'ensemble de ses oeuvres numérisées disponible dans la bibliothèque Cujas ou sur des sites partenaires.

Une deuxième forme d'éditorialisation des collections sera d'écrire des billets de blog, des parcours thématiques, écrite par des spécialistes d'une question ou d'une thématique pour d'autre s'intéressant au sujet. Cela pourrait également comme le suggérait l'un des enseignats-chercheurs une porte d'entrée pour les jeunes chercheurs pour développer leur carrière. L'ancienne responsable de la bibliothèque numérique rappelle l'intérêt d'associer directement les étudiants avancés et les chercheurs à cette démarche de valorisation. L'idée serait de permettre à des chercheurs ou des étudiants d'écrire des articles des billets de blog sur un ou des documents étudiés pour leur étude provenant de la bibliothèque Cujas à le replacer dans un contexte intellectuel, redonnant \enquote{vie} aux documents numérisés. Laurent Pfister pense que ce serait un excellent moyen de valoriser la production scientifique notamment pour des jeunes chercheurs.

En combinant la flexibilité technique du format Markdown et la rigueur des guides d'expositions de la BIS, la bibliothèque numérique de Cujas ne sera plus uniquement un outil d'exposition de ses documents numérisés, mais deviendrait une véritable plateforme de médiation culturelle juridique. POur que ces services soient utilisés par les usagers, il est nécessaire pour les experts métiers de les accompagner dans leur prise en main. Il est important alors que les experts métiers eux-mêmes soient capables de les utiliser amenant le service dédié de la bibliothèque à les accomapgner dans la conduite du changement.